\begin{theorem}[Fundamental theorem of algebraic graph theory]
    \thmlab{{f:t:algebraic:graph}}%
    %
    Let $\Graph= \Graph(V,E)$ be an undirected (multi)graph with
    maximum degree $d$ and with $n$ vertices.  Let
    $\lambda_1\geq\lambda_2\geq \cdots \geq \lambda_n$ be the
    eigenvalues of $\Mtrx(\Graph)$ and the corresponding orthonormal
    eigenvectors are $e_1,\ldots, e_n$. The following holds.
    \begin{compactenumi}
        \item If $\Graph$ is connected then $\lambda_2 < \lambda_1$.

        \item For $i=1,\ldots, n$, we have
        $\cardin{\lambda_i} \leq d$.

        \item $d$ is an eigenvalue if and only if $\Graph$ is regular.

        \item If $\Graph$ is $d$-regular then the eigenvalue
        $\lambda_1=d$ has the eigenvector $e_1 = \frac{1}{\sqrt{n}}
        (1,1,1, \ldots, 1)$.

        \item The graph $\Graph$ is bipartite if and only if for every
        eigenvalue $\lambda$ there is an eigenvalue $-\lambda$ of the
        same multiplicity.

        \item Suppose that $\Graph$ is connected. Then $\Graph$ is bipartite if
        and only if $-\lambda_1$ is an eigenvalue.

        \item If $\Graph$ is $d$-regular and bipartite, then $\lambda_n =
        d$ and $e_n=\frac{1}{\sqrt{n}} ( 1,1, \ldots, 1, -1, \ldots,
        -1)$, where there are equal numbers of $1$s and $-1$s in
        $e_n$.
    \end{compactenumi}
\end{theorem}
